\section{Introduction}
Reference-guided dense medical segmentation can use frozen matching evidence to propose positive and negative prompts for a promptable segmenter. AdaPromptSet frames this interface as budget-constrained selection among complete structured prompt sets. The intended adaptive question is deliberately narrow: under a fixed matcher, fixed segmenter, fixed candidate generator, and explicit SAM-call budget, does query-specific prompt-set choice have enough measurable headroom to justify a learned selector? This framing separates prompt construction from model adaptation. It also makes failure interpretable: if the best validation-selected fixed action is already close to the per-query oracle, then training a selector would be a narrative convenience rather than an evidence-backed method.

Current MoNuSAC evidence records \OracleAnalysisDecision{} and keeps policy training \OracleAnalysisPolicyTraining{}. The main development pilot used ten validation episodes and twelve valid v2 actions. The best fixed action, \OracleAnalysisBestFixed{}, achieved mean Dice \OracleAnalysisBestFixedDice{}, while the per-query oracle achieved \OracleAnalysisOracleDice{}. The measured Dice gap, \OracleAnalysisGap{}, is below the predeclared minimum needed to continue policy learning. This result does not prove that all prompt selection is unnecessary; it says that the present observable evidence, action bank, and MoNuSAC split do not support a stronger adaptive-policy claim.

The contribution of this paper is therefore procedural and diagnostic rather than promotional. First, it gives a reproducible gate sequence for deciding when prompt-set selection should stop before policy training. Second, it reports an action-oracle analysis showing how much benefit remains after selecting the strongest fixed rule. Third, it records why the present failure should not be hidden: candidate coverage is limited, higher-call grouping did not show a clear gain, and the statistical unit count is too small for a formal positive conclusion. These points are useful precisely because they constrain the next experiment instead of expanding the method after seeing the outcome.
